% 冻结 v6；全局符号对应 GitHub c37feaf3 的“梳理/论文初稿.tex”。
% 本文件可直接替换总稿的“问题三的模型与求解”一节。
\section{问题三的模型与求解}
\label{sec:q3-v6}

\subsection{从单源定位到搜索与清除协同}

前两问讨论了测向信息如何约束源位置，以及改变检测位置对定位效果的影响。第三问还需解决数量未知时的完整搜索：机器狗既要处理已发现的源，也要排查尚未发现源的频道，而两类任务需要共同使用一条移动路线。若每发现一个源便单独完成定位和清除，再返回固定站点继续扫描，可能产生较多往返；若先扫描完整个区域再处理目标，又难以及时利用新观测调整行程。因此，本文采用基于位置可行区域的搜索—定位—清除协同策略，在执行过程中持续更新信息和后续访问计划。

Vander Hook等在主动纯方位定位中同时考虑移动时间与测量时间，说明检测点的几何效果需要与到达该点的费用共同评价\cite{VanderHook2014}。结合本题计费规则，记总移动距离为$D_{\mathrm{move}}$，无线电检测、实际换频、成功清除和失败清除次数分别为$n_{\mathrm{meas}}$、$n_{\mathrm{sw}}$、$n_{\mathrm{ok}}$和$n_{\mathrm{fail}}$，则总虚拟时间为
\begin{equation}
 T_{\mathrm{total}}=\frac{D_{\mathrm{move}}}{5}+5n_{\mathrm{meas}}+n_{\mathrm{sw}}
                     +5n_{\mathrm{ok}}+3n_{\mathrm{fail}}.
 \label{eq:q3-time}
\end{equation}
式中距离以米计、时间以秒计。清除成功的5秒包括光学定位3秒与清除2秒；清除操作不改变测向机的当前频道。无信号检测和确认不存在剩余源的动作也计入费用，任务结束后无需返回原点。本文在满足完整清除要求的条件下，利用预计时间比较候选路线与局部动作。

沿用全局符号，以$K$表示综合观测信息得到的位置可行区域，以$c(K)$、$\rho(K)$和$d_{\max}(q;K)$分别表示最小包围圆圆心、半径和候选行动点到区域的最远距离。多源情形以频道$j$区分目标，用$K_{j,t}$表示第$t$次信息更新后的区域，$p_t$表示机器狗当前位置。记已发现频道集合为$\mathcal D_t$，已清除频道集合为$\mathcal C_t$；尚未发现且未被完整排除的频道组成$\mathcal J_t$。这些量均由实际反馈更新，未知真实位置$g_j$与实际源数$N$不作为在线决策输入。

\begin{figure}[htbp]
 \centering
 \includegraphics[width=\linewidth]{q3-flow.pdf}
 \caption{搜索、定位与清除协同策略的执行流程}
 \label{fig:q3-flow}
\end{figure}

图\ref{fig:q3-flow}给出策略的整体结构。位置可行区域负责判断如何继续定位和能否可靠清除，逐频道覆盖负责确认没有漏源，任务排序与局部费用比较负责减少行动时间。路线和费用估计随反馈更新，而完成判定始终依据实际取得的观测。

\subsection{由观测约束到可靠清除}

\subsubsection{融合测向、接收与清除反馈}

一次正常测向将源位置限制在正向角域$W_i$内。依据第一问的半平面表示，叠加新的测向信息即可缩小原区域。对全向源，收到信号还说明源距不超过1500米；强信号进一步说明源距不超过$r_{\mathrm{strong}}$。在检测点$s_i$获取反馈后，可分别按下列约束更新：
\begin{equation}
\begin{aligned}
 K_{j,t+1}&=K_{j,t}\cap W_i\cap B(s_i,1500), &&\text{正常示向度},\\
 K_{j,t+1}&=K_{j,t}\cap B(s_i,r_{\mathrm{strong}}), &&\text{强信号},\\
 K_{j,t+1}&=K_{j,t}\setminus B(s_i,1000), &&\text{无信号}.
\end{aligned}
\label{eq:q3-region-update}
\end{equation}
其中$B(s,r)$为以$s$为圆心、$r$为半径的闭圆盘；未知频道在首次约束前的位置范围取为$\Omega$。这些更新保留所有与所用约束相容的位置，不把区域中心视为真实源位置。

同一源的接收半径$R_j$虽然未知，但在检测过程中固定。因此，若$s_+$有信号而$s_-$无信号，则真实位置满足
\begin{equation}
 \|g_j-s_+\|\le R_j<\|g_j-s_-\|,
 \qquad
 2(s_--s_+)^{\mathsf T}g_j<\|s_-\|^2-\|s_+\|^2.
 \label{eq:q3-coupling}
\end{equation}
这给出额外的半平面约束，能够利用正、负观测之间的联系进一步收缩区域。

几何计算采用保守外包区域。读数保留至$0.01^\circ$时，将题设误差界$\alpha$扩大至$1.0051^\circ$对应的弧度，以包含舍入和微小数值裕量；圆域采用适当的多边形外包。负观测区域排除后，可取凸包供后续测向几何计算使用，该处理可能保留多余位置，但不应删除真实位置。光学试清失败产生的孔洞或分裂区域则另行保留，避免后续再次把已排除部分当作新的试清收益。

\subsubsection{根据区域尺度选择观测与清除位置}

由第一问，位置$q$能够保证进入源的光学清除范围，当且仅当
\begin{equation}
 d_{\max}(q;K_{j,t})=\max_{x\in K_{j,t}}\|q-x\|\le r_{\mathrm{opt}}.
 \label{eq:q3-safe}
\end{equation}
实现中取安全距离$r_{\mathrm{safe}}=19.5\,\mathrm m$。当$\rho(K_{j,t})\le r_{\mathrm{safe}}$时，从$c(K_{j,t})$朝$p_t$的线段上选择满足式\eqref{eq:q3-safe}且以$r_{\mathrm{safe}}$替代$r_{\mathrm{opt}}$的最靠近$p_t$的位置，随后执行清除。这样可以在保留安全裕量的同时减少最后一段移动。这里仅在线段上选点，不将其称为整个可清除区域上的全局最近点。

若区域尚大，则继续选择有助于交会的检测位置。记区域最长方向的单位法向量为$v_{j,t}$。当$\rho(K_{j,t})>120\,\mathrm m$时，构造侧向候选点
\begin{equation}
 q_\pm=c(K_{j,t})\pm
 \min\{150,\,0.25\rho(K_{j,t})\}\,v_{j,t};
 \label{eq:q3-next-point}
\end{equation}
区域较小时优先考虑中心。候选点须通过接收条件检验，再从中选择移动较近者。例如，若整个区域距$q$不超过1000米，则必能接收到信号；也可利用首次正观测和接收半径下界构造充分的接收条件。由于同地点误差固定，对重复检测位置改用满足接收条件的侧移点，避免依靠原地重复测量获得虚假的精度改善。

\subsection{以逐频道覆盖保证搜索完整}

处理已发现目标只能说明$\mathcal D_t\setminus\mathcal C_t$是否为空，不能排除未知频道仍有遗漏。对尚未发现的频道$j$，记实际取得无信号反馈的位置集合为$\mathcal P_{j,t}^{-}$。由于全向源接收半径至少为1000米，这些观测对应的尚未排除区域为
\begin{equation}
 \mathcal U_{j,t}=\Omega\setminus
 \bigcup_{s\in\mathcal P_{j,t}^{-}}B(s,1000).
 \label{eq:q3-unsearched}
\end{equation}
若$\mathcal U_{j,t}=\varnothing$，则该频道不可能仍有未发现源。覆盖记录必须逐频道保存：经过某点但没有检测该频道，或只检测了其他频道，都不能计入$\mathcal P_{j,t}^{-}$。

为了将连续覆盖判定落到计算中，把包围$\Omega$的正方形划分为边长28.125米的网格，保留所有与$\Omega$相交的单元，包括边界单元。设单元$Q_\ell$的四个角点集合为$V_\ell$。仅当存在该频道的一次实际负观测$s$满足
\begin{equation}
 \max_{v\in V_\ell}\|v-s\|\le999.99
 \label{eq:q3-cell-cover}
\end{equation}
时，才将整个单元标记为已排除。由于圆盘是凸集，四角均在同一排除圆内即可保证整个单元被覆盖。所有保留单元均通过检查，是式\eqref{eq:q3-unsearched}为空的充分条件；使用略小于1000米的半径留出数值裕量。这一条件较保守，但避免了仅检查网格中心导致的漏判。

因此，正常完成任务须满足
\begin{equation}
 \mathcal D_t\setminus\mathcal C_t=\varnothing,
 \qquad
 \mathcal U_{j,t}=\varnothing\quad\text{对所有未发现频道 }j,
 \label{eq:q3-stop}
\end{equation}
或已有$\#\mathcal C_t=16$，达到题设源数上限。达到上限之前，连续若干次无信号或暂时没有待处理目标，均不能代替这一完成条件。

\subsection{联合安排扫描与目标处理任务}

Engin和Isler的多目标主动定位研究将平台运动与多个目标的信息需求共同考虑\cite{Engin2020}。这一研究提示，同一段移动可以服务多个目标。结合本题源数未知、逐频道付费检测的特点，本文将必要的覆盖扫描与已发现源的处理共同纳入访问计划；采用显式几何和路线评价完成这一安排。

\subsubsection{构造覆盖候选并调整站点}

初始时在原点扫描20个频道，再设置长度250米的辅助基线，根据已发现目标方向安排侧向位置，并选择值得补测的频道。在此基础上，以原点和半径1200米圆周上的六个候选点构造搜索布局，分别考察$0^\circ,10^\circ,\ldots,50^\circ$的整体旋转。已发现目标的预计处理位置也可作为共享扫描位置。

每次规划均利用实际覆盖记录删除冗余扫描任务，并调整剩余站点。调整时保留该站点必须覆盖的网格角点约束，在满足覆盖条件的范围内缩短其与相邻任务之间的路程。因此，六点布局只是构造可行候选的起点，执行中不要求按固定七站顺序巡游。计划中的覆盖仅用于安排未来动作；完成判定仍须等待相应频道的实际负观测。

\subsubsection{用处理任务的进入与离开位置评价路线}

扫描任务在站点执行，目标处理任务则可能先到新检测点，再到清除点。为反映这一区别，记任务$b$的预计离开位置为$e_b$：扫描任务取其站点，目标处理任务以$c(K_{j,t})$作近似。从前一任务$a$的预计离开位置出发，记任务$b$的进入位置为$q_b(e_a)$，对目标任务取可行的清除位置或下一测向位置。定义有向路程代理
\begin{equation}
 \widehat d_{ab}=\|q_b(e_a)-e_a\|+\|e_b-q_b(e_a)\|.
 \label{eq:q3-route-distance}
\end{equation}
该代理包含进入处理过程的路程，通常不具有对称性。以当前位置为序列起点，对候选任务序列$\pi$计算
\begin{equation}
 \widehat T_{\mathrm{plan}}(\pi)
 =\frac{1}{5}\sum_{(a,b)\in\pi}\widehat d_{ab}
   +6\,\#\mathcal J_t\,m_{\mathrm{scan}}(\pi),
 \label{eq:q3-plan-score}
\end{equation}
其中$m_{\mathrm{scan}}(\pi)$为序列中的扫描站数，6秒对应一次测量及一次换频的保守估价。通过多起点近邻构造和路径片段反转改进比较访问序列，仅执行首项任务，取得反馈后重新规划。式\eqref{eq:q3-plan-score}是候选方案的时间代理，不是整个剩余任务时间的精确价值函数；实际费用按式\eqref{eq:q3-time}累计。

较长移动途中，在行程的四分之一、二分之一和四分之三处检查已发现源的补测机会。若候选位置与已有检测点有足够间距，且代表性方向下预计包围圆半径明显减小，则选择相应频道检测。沿途补测采用几何收益筛选，并逐频道付费，不将一次停点视为同时免费观测所有目标。

\subsection{继续测向与提前试清的局部时间比较}

可靠清除要求整个可行区域落入同一清除圆。当区域尚未满足该条件时，直接试清仍可能成功，但也可能产生失败费用。因此，仅比较“一次测量5秒”和“一次试清失败3秒”并不足以作出决策，还应计入到达行动点的移动以及失败后的继续定位时间。

当$r_{\mathrm{safe}}<\rho(K_{j,t})\le100\,\mathrm m$时，从当前位置、区域中心及少量邻近几何位置生成试清候选。剔除与既有试清位置过近，或只能覆盖极少可行面积的候选。取区域内的少量代表位置，并与$-\alpha,0,\alpha$三种误差组合，形成有限情形集合$\Xi_{j,t}$。对情形$\xi=(x,e)$和候选试清位置$q$，估计当前源完成处理的时间为
\begin{equation}
 \widehat T_{\mathrm{try}}(q;\xi)
 =\frac{\|q-p_t\|}{5}+
 \begin{cases}
  5,&\|x-q\|\le r_{\mathrm{opt}},\\
  3+\widehat T_{\mathrm{cont}}(K_{j,t}^{-}(q),q;\xi),&\|x-q\|>r_{\mathrm{opt}},
 \end{cases}
 \label{eq:q3-try-time}
\end{equation}
其中$K_{j,t}^{-}(q)$为试清失败后的剩余区域，$\widehat T_{\mathrm{cont}}$由后续几何测向与清除步骤有限步试算得到。继续无线电定位的时间$\widehat T_{\mathrm{radio}}(\xi)$也按同一局部完成口径计算。试算最多推进五次普通测向，未完成的情形追加1000秒惩罚，避免将截断前缀误作已完成的费用。

对操作$a$，以平均费用与较差情形费用的加权值作为评分：
\begin{equation}
 J_{\mathrm{loc}}(a)=
 0.8\frac{1}{\#\Xi_{j,t}}\sum_{\xi\in\Xi_{j,t}}\widehat T_a(\xi)
 +0.2\max_{\xi\in\Xi_{j,t}}\widehat T_a(\xi).
 \label{eq:q3-local-score}
\end{equation}
当最优候选试清评分比继续测向低至少1秒时，执行试清；否则继续测向，每个源最多提前试清两次。该评分使用有限代表情形，是降低计算量的启发式估价，不解释为真实后验期望或连续状态上的最坏时间保证。

若实际试清失败，由题设可知$\|g_j-q\|>r_{\mathrm{opt}}$，因而可排除相应圆域。实现中从已有区域去除半径19.999米圆的内接多边形，保持外包性质，并保存可能出现的非凸或分裂结构。随后根据新的区域继续比较动作，失败信息不会丢失。

\begin{figure}[htbp]
 \centering
 \includegraphics[width=\linewidth]{q3-optical.pdf}
 \caption{同一官方演练中频道3的试清与后续定位。图(a)蓝色为试清前区域，绿色为失败后保留部分，橙色圆为20米试清范围；图(b)灰色为失败后区域，蓝色为再测向后的区域，虚线圆为19.5米安全清除范围。局部坐标$x',y'$由原坐标平移、旋转得到，图形未使用真实源坐标。}
 \label{fig:q3-optical}
\end{figure}

图\ref{fig:q3-optical}展示一个实际执行片段。试清前区域半径为87.43米，试清与继续测向的局部评分分别为24.90秒和29.59秒，因此选择试清。实际失败后，区域分为两部分，包围圆半径仍约为87.43米；这说明排除了部分面积，并不一定立即降低最坏位置尺度。随后在试清位置测向，区域半径缩小至18.28米，再移动到满足$d_{\max}=19.5$米的位置成功清除。这里的局部评分用于说明当时的选择，不作为已实现节省4.69秒的反事实结论。

\begin{table}[htbp]
 \centering\small
 \caption{算法主要几何与局部决策参数}
 \label{tab:q3-parameters}
 \renewcommand{\arraystretch}{1.16}
 \begin{tabularx}{\linewidth}{lclY}
 \toprule
 参数 & 取值 & 单位 & 作用\\
 \midrule
 安全清除距离$r_{\mathrm{safe}}$ & 19.5 & m & 在物理20米范围内留出裕量\\
 辅助观测基线 & 250 & m & 改善起点观测的交会几何\\
 初始外围候选半径 & 1200 & m & 构造覆盖扫描布局\\
 外围布局旋转数 & 6 & -- & 比较不同覆盖与访问组合\\
 覆盖网格边长 & 28.125 & m & 对整个单元进行保守排除检查\\
 试清区域半径上限 & 100 & m & 限制提前试清的适用范围\\
 每源提前试清上限 & 2 & 次 & 限制局部试清开销\\
 试清评分改善门槛 & 1 & s & 仅采用预计费用更低的候选\\
 \bottomrule
 \end{tabularx}
\end{table}

表\ref{tab:q3-parameters}列出执行参数。这些值属于冻结实现的选择，不由文献的理论保证推出。全程的完成依据仍为式\eqref{eq:q3-stop}及源数上限，而不是局部评分大小或计算循环次数。

\subsection{官方演练结果与策略分析}

\subsubsection{整体完成情况与时间组成}

采用本文策略的两批官方演练记录检验执行效果。设第$b$场实际源数为$N_b$、总虚拟时间为$T_b$，以场为单位计算
\begin{equation}
 \overline\tau=\frac{1}{B}\sum_{b=1}^{B}\frac{T_b}{N_b}.
 \label{eq:q3-evaluation}
\end{equation}
这一指标与所有场总时间除以总源数的加权平均不同。各场源数仅由演练结果用于事后评价；两批场景分别统计，不把批次差异解释为策略改进。

\begin{table}[htbp]
 \centering\small
 \caption{两批官方演练的完整清除与耗时结果}
 \label{tab:q3-results}
 \renewcommand{\arraystretch}{1.2}
 \begin{tabularx}{\linewidth}{Zrrrrrr}
 \toprule
 批次 & 全清场数 & 源总数 & 均值 & 中位数 & 标准差 & 最大值\\
 \midrule
 50场 & 50/50 & 645 & 231.71 & 227.17 & 31.20 & 303.41\\
 30场 & 30/30 & 395 & 226.34 & 213.86 & 31.03 & 293.51\\
 \bottomrule
 \multicolumn{7}{l}{\footnotesize 注：后四列单位均为s/源；全部为官方演练，非正式测试成绩。}\\
 \end{tabularx}
\end{table}

两批记录均完成各场全部干扰源的清除，场均每源耗时分别为231.71秒和226.34秒。图\ref{fig:q3-performance}(a)保留每场观测值，并用箱线图显示离散程度，表明不同场景下仍有明显的费用差异。源数量、空间位置和接收条件均可能影响行程，不能仅以一次较快演练代表整体效率。

\begin{figure}[htbp]
 \centering
 \includegraphics[width=\linewidth]{q3-performance.pdf}
 \caption{本文策略的演练耗时分布及组成。图(a)每点对应一场，菱形为场均值，箱体为四分位区间，中线为中位数，须线采用1.5倍四分位距规则；图(b)各项均先除以本场实际源数，再对场次平均。}
 \label{fig:q3-performance}
\end{figure}

\begin{table}[htbp]
 \centering\small
 \caption{场均每源时间的分项构成（s/源）}
 \label{tab:q3-costs}
 \renewcommand{\arraystretch}{1.15}
 \begin{tabularx}{.86\linewidth}{Yrr}
 \toprule
 时间项目 & 50场批次 & 30场批次\\
 \midrule
 移动 & 170.72 & 165.57\\
 无线电检测 & 46.70 & 46.50\\
 换频 & 8.86 & 8.84\\
 成功清除 & 5.00 & 5.00\\
 失败清除 & 0.44 & 0.43\\
 \midrule
 总计 & 231.71 & 226.34\\
 \bottomrule
 \multicolumn{3}{l}{\footnotesize 注：分项之和与总计的微小差异来自分别舍入。}\\
 \end{tabularx}
\end{table}

由表\ref{tab:q3-costs}与图\ref{fig:q3-performance}(b)，移动约占50场批次总时间的73.68\%，是主要费用来源，支持将覆盖与目标处理放在同一路线中考虑的建模动机。该批次无线电检测约为9.34次/源，其中首次发现及此前搜索约占7.14次，发现后定位约占2.20次。可见，总检测次数还包含空频道排除与首次发现前的搜索，不能全部归为单源交会定位。

两批分别发生89次和57次提前清除失败，直接费用约为0.44秒/源和0.43秒/源，均已包含在总时间中。这些是动作失败，而非场景未全清。现有记录能够检验完整策略的实际表现，但未单独识别每个模块的因果收益，因此不据此宣称提前试清或联合路由各自节省了确定比例的时间。

\subsubsection{代表场景中的路线与完成证据}

从50场批次按每源耗时升序选取第25场，即较低中位场，展示完整执行过程，而不选择最快场景。该场共清除13个源，总移动距离为10396.34米，无线电检测135次，实际换频127次，清除失败1次。由式\eqref{eq:q3-time}可得
\begin{equation}
 T_{\mathrm{total}}=\frac{10396.34}{5}+5\times135+127+5\times13+3
 \approx2949.27\,\mathrm s,
\end{equation}
对应226.87秒/源，与官方记录一致。

\begin{figure}[htbp]
 \centering
 \includegraphics[width=\linewidth]{q3-route-coverage.pdf}
 \caption{较低中位演练场景的实际路线与覆盖证据。图(a)连线依照动作顺序连接，星号为实际成功清除位置而非源的真实坐标；图(b)为频道6的9个实际无信号检测点及其保证排除圆，覆盖图仅展示任务圆内部分。}
 \label{fig:q3-route}
\end{figure}

图\ref{fig:q3-route}(a)中，检测与清除位置交错出现，说明搜索任务实际穿插在目标处理过程中。图\ref{fig:q3-route}(b)展示一个不存在源的频道所积累的负观测：9个排除圆共同覆盖任务区域。对包含边界的13104个网格单元重新检查，均满足式\eqref{eq:q3-cell-cover}。本场其余6个未发现频道也分别通过覆盖检查，已发现的13个频道均有成功清除反馈，因此具备正常终止的依据。

综合两批演练与代表场景，所建策略将位置约束、搜索完整性和行动费用连接为可执行过程，在已测试场景中完成全部清除。其局部有限情形估价及路线代理仍具有近似性，样本全清不等同于对所有合法场景的软件有限步完成证明，也不构成全局时间最优性结论。
